\documentclass[conference]{IEEEtran}
\IEEEoverridecommandlockouts

\usepackage{cite}
\usepackage{amsmath,amssymb,amsfonts}
\usepackage{algorithmic}
\usepackage{graphicx}
\usepackage{textcomp}
\usepackage{xcolor}
\usepackage{hyperref}

\hypersetup{hidelinks}

\begin{document}

\title{HW-Triage-Bench: A Benchmark for Hardware Test Failure-to-Report Automation}

\author{\IEEEauthorblockN{Piyush Jagadish Bag}
\IEEEauthorblockA{\textit{Test Automation Engineer, Platform Hardware}\\
Santa Clara, CA, USA\\
piyushbag4@gmail.com}}

\maketitle

\begin{abstract}
Hardware validation engineers spend more time debugging and documenting failures than running tests.
Industry surveys report debugging as the dominant verification activity, and recent academic work notes
that automated analysis of engineering validation logs still lacks a standard, reproducible approach.
Commercial electronic design automation flows automate simulation and formal verification but typically
stop before the post-test step: assembling failure evidence, locating affected components in schematics,
and producing a reviewer-ready report remains manual across large program portfolios.

We introduce \textbf{HW-Triage-Bench}, an open benchmark and reference pipeline for the
\textbf{failure-to-report} path. Given heterogeneous validation logs and schematic PDFs, systems must
classify failure signatures, retrieve relevant design context, and emit structured diagnostic summaries
suitable for design review. We evaluate rule-based, retrieval-augmented, and agentic baselines on
\textit{N} anonymized projects spanning \textit{M} failure classes, reporting precision and recall on log
classification and human-judged report usefulness.

We release benchmark tasks, evaluation harness code, and a non-proprietary \textbf{OpenTestTriage}
reference module integrated with open lab stacks. HW-Triage-Bench establishes a public baseline for a
gap the field describes as open but rarely measures.
\end{abstract}

\begin{IEEEkeywords}
hardware validation, test log analysis, post-silicon debug, benchmark, design review automation
\end{IEEEkeywords}

\section{Introduction}
\label{sec:intro}
% TODO: Wilson 2024 debugging time; Kanagal Berkeley 2025 EDA log gap; gap = failure-to-report

\subsection{Contributions}
\begin{itemize}
    \item \textbf{HW-Triage-Bench}, the first open benchmark for validation failure-to-report automation.
    \item A reference \textbf{OpenTestTriage} pipeline: log classification, schematic localization, report assembly.
    \item Baseline evaluation across rule-based, retrieval-augmented, and agentic matchers.
    \item Open harness aligned with OCP diagnostic output schemas.
\end{itemize}

\section{Problem Statement}
\label{sec:problem}

Modern hardware validation produces far more failure evidence than any single review cycle can absorb manually.
Functional verification surveys consistently rank \emph{debugging} as the largest share of engineer time~\cite{wilson2024verification}.
Yet the bottleneck is not only root-cause isolation in silicon or on the bench: after tests run, teams must still transform raw logs, database queries, and schematic PDFs into a coherent, reviewer-ready report.
Recent academic work on EDA log analysis notes that this post-test documentation step still lacks a standard automated approach~\cite{kanagal2025eda}.

We call this the \textbf{failure-to-report gap}.
Commercial EDA and lab-automation flows automate simulation, pattern execution, and instrument capture, but typically stop before three operations that remain manual in practice:
\begin{enumerate}
    \item \textbf{Log reconciliation} --- sorting and classifying hundreds to thousands of heterogeneous validation logs per hardware program, often with inconsistent severity labels and duplicate signatures across runs.
    \item \textbf{Design-context localization} --- mapping each classified failure to the correct components and regions in multi-page schematic PDFs so reviewers know \emph{where} to look.
    \item \textbf{Report assembly} --- merging classified failures, localized schematic context, and program metadata into a single structured document suitable for cross-functional design review.
\end{enumerate}

At portfolio scale, manual assembly dominates calendar time.
Industry practitioners report on the order of one engineer-hour per program per review cycle when reports are built by hand, repeated across hundreds of active programs, and dominated by search-and-copy rather than analysis~\cite{kanagal2025eda}.
The cost is super-linear: each new product generation adds log formats, codenames, and schematic revisions, while review deadlines do not scale accordingly.

Existing partial solutions address slices of the pipeline but not the end-to-end path.
Pre-silicon verification tools focus on coverage, formal properties, and simulation closure~\cite{seshia2010postsi}.
Knowledge-graph and LLM methods for \emph{textual} test reports target compliance or aerospace board data~\cite{thales2024llmkg}, not the recurring hardware-design review workflow that joins validation logs with schematic artifacts.
Open diagnostic schemas such as OCP diag-core standardize \emph{what} a test emits~\cite{ocpdiagcore}, but not how failures become reviewer-ready documentation across inconsistent internal metadata.

\textbf{Problem definition.}
Given a hardware program's validation log corpus and associated schematic PDFs, produce a structured diagnostic report that (i)~classifies failure signatures with reproducible precision/recall, (ii)~links each signature to schematic context, and (iii)~reduces human assembly time from hours to minutes without exposing proprietary design data.
No public benchmark currently measures this full path.
HW-Triage-Bench (Section~\ref{sec:benchmark}) and the OpenTestTriage reference pipeline (Section~\ref{sec:pipeline}) are designed to close that measurement and reproducibility gap while keeping the core problem anchored in post-test validation documentation for schematic-centric design review.

\section{Related Work}
\label{sec:related}
% TODO: EDA verification automation, LLM log analysis, Thales test-data KG (arXiv:2408.01700)

\section{HW-Triage-Bench}
\label{sec:benchmark}
% TODO: task definition, dataset construction, anonymization protocol, metrics

\section{OpenTestTriage Reference Pipeline}
\label{sec:pipeline}
% TODO: tiered fuzzy classification, PDF localization, report template — no proprietary DB names

\section{Evaluation}
\label{sec:eval}
% TODO: baselines, F1, report rubric, latency

\section{Threats to Validity}
\label{sec:threats}
% TODO: synthetic vs anonymized real logs; generalization across vendors

\section{Conclusion}
\label{sec:conclusion}
% TODO: future work: agentic extensions, integration with OpenTAP/labgrid

\section*{Data Availability}
Benchmark artifacts and evaluation code: \url{https://github.com/piyushbag/open-test-triage}.
Anonymized dataset DOI: Zenodo TBD.

\bibliographystyle{IEEEtran}
\bibliography{references}

\end{document}
